\documentclass[aspectratio=169]{beamer}
\usetheme{metropolis}
\usepackage{booktabs}
\usepackage{amsmath,amssymb}
\usepackage{listings}

\lstset{
  basicstyle=\ttfamily\scriptsize,
  keywordstyle=\color{blue},
  commentstyle=\color{gray},
  breaklines=true,
  frame=none,
}

\title{Bispectrum Library: From Finite Groups to the Disk}
\subtitle{Implementing Mataigne et al.\ (NeurIPS 2024) \& Myers \& Miolane (arXiv 2025)}
\author{Johan Mathe}
\date{\today}

\begin{document}

\maketitle

% ============================================================
\begin{frame}{Goal}
  Reimplement the selective $G$-bispectrum from Mataigne et al.\ (NeurIPS 2024)
  in a standalone, clean PyTorch library.

  \bigskip

  Three group/domain families implemented:
  \begin{itemize}
    \item \textbf{CnonCn} --- cyclic groups $C_n$ acting on $\mathbb{Z}/n\mathbb{Z}$ \hfill\textit{Mataigne et al.}
    \item \textbf{DnonDn} --- dihedral groups $D_n$ acting on $D_n$ \hfill\textit{Mataigne et al.}
    \item \textbf{SO2onD2} --- $\mathrm{SO}(2)$ acting on the unit disk $\mathbb{D}^2$ \hfill\textit{Myers \& Miolane}
  \end{itemize}

  \bigskip

  All three include \texttt{forward()} (bispectrum) and \texttt{invert()} (signal recovery).

  \smallskip

  Plus \textbf{SO3onS2} ($\mathrm{SO}(3)$ on $S^2$) --- forward only, no inversion.
\end{frame}

% ============================================================
\begin{frame}{Improvements over g-invariance repo}
  \begin{table}
    \scriptsize
    \begin{tabular}{lll}
      \toprule
      \textbf{Aspect} & \textbf{g-invariance} & \textbf{bispectrum (new)} \\
      \midrule
      Dependencies & escnn, einops, scipy & Pure PyTorch \\
      CG computation & \texttt{scipy.linalg.schur} & \texttt{torch.linalg.eig} \\
      Group size & $n \geq 8$ enforced & All $n > 2$ \\
      $n_3$ for odd $n$ & Bug: gives 0 for $n=3$ & Fixed \\
      Inversion & Not implemented & Algorithms 2 \& 4 \\
      Interface & escnn \texttt{GeometricTensor} & Raw \texttt{torch.Tensor} \\
      Normalization & Baked in & None (raw coefficients) \\
      Group DFT & $O(n^2)$ explicit & $O(n \log n)$ via FFT \\
      Tests & \textbf{None} & \textbf{180 tests} \\
      \bottomrule
    \end{tabular}
  \end{table}
\end{frame}

% ============================================================
\begin{frame}{Bugs found in g-invariance (1/2)}
  \textbf{Bug 1: $n_3$ calculation breaks for small odd $n$}

  \smallskip

  g-invariance code: \texttt{n3 = n2 - 1} where $n_2 = \lfloor(n{-}1)/2\rfloor$.

  For $n=3$: $n_2 = 1$, so $n_3 = 0$ --- meaning \emph{zero}
  $\beta_{\rho_1,\rho_k}$ coefficients are computed.

  \smallskip

  \textbf{Impact}: $\beta_{\rho_1,\rho_1}$ is needed to recover $\mathcal{F}(\rho_{01})$
  (Appendix~E of the paper proves $\rho_{01} \in \rho_1 \otimes \rho_1$).
  Without it, the selective bispectrum for $D_3$ is \emph{incomplete} ---
  the 1D irrep $\rho_{01}$ cannot be recovered during inversion.

  \smallskip

  \textbf{How discovered}: Writing parametrized tests over $n \in \{3, 4, 5, 7, 8\}$.
  $D_3$ inversion failed because \texttt{output\_size} was only $1 + 4 = 5$
  (missing the $16 \times n_3$ block entirely).
  The g-invariance repo never tested $n < 8$ due to bug~2.

  \smallskip

  \textbf{Fix}: \texttt{n3 = n2 if n \% 2 == 0 else max(n2 - 1, 1)}.
\end{frame}

% ============================================================
\begin{frame}{Bugs found in g-invariance (2/2)}
  \textbf{Bug 2: Artificial $n \geq 8$ restriction}

  \smallskip

  \texttt{pooling.py} line 390: \texttt{if n < 8: raise ValueError}.

  \smallskip

  This came from escnn integration constraints, not from the math.
  $D_n$ bispectrum is well-defined for all $n > 2$. Dropping this
  restriction immediately exposed Bug~1.

  \bigskip

  \textbf{Bug 3: Paper typo in Algorithm 4} (next slide).

  \bigskip

  \textbf{Note}: The g-invariance repo had \textbf{zero unit tests},
  so these bugs were never caught. Bug~1 is silent --- it doesn't crash,
  it just computes an incomplete bispectrum.
\end{frame}

% ============================================================
\begin{frame}{Paper typo in Algorithm 4 (bispectrum inversion on $D_n$)}
  Algorithm 4, line 6 computes the block-diagonal Fourier coefficients via:
  \[
    \bigoplus_{\rho \in \rho_1 \otimes \rho_{k-1}} \mathcal{F}(\Theta)_\rho
    = \Bigl(
      C_{\rho_1,\rho_{k-1}}^\dagger
      \bigl[\mathcal{F}_{\rho_1} \otimes \mathcal{F}_{\rho_{k-1}}\bigr]^{-1}
      \beta_{\rho_1,\rho_{\mathbf{2}}}
      \, C_{\rho_1,\rho_{k-1}}
    \Bigr)^\dagger
  \]

  The subscript on $\beta$ is printed as $\rho_2$ (fixed), but should be
  $\rho_{k-1}$ (varying with the loop index $k$). The CG subscripts
  $C_{\rho_1,\rho_{k-1}}$ are correct.

  \bigskip

  \textbf{How discovered}: When implementing inversion for $n=8$ (which has
  $\lfloor\frac{7}{2}\rfloor = 3$ 2D irreps), the loop needs
  $\beta_{\rho_1,\rho_1}$ then $\beta_{\rho_1,\rho_2}$ to recover
  $\mathcal{F}(\rho_2)$ then $\mathcal{F}(\rho_3)$.
  Using a fixed $\beta_{\rho_1,\rho_2}$ for both iterations gives wrong
  Fourier coefficients --- the Frobenius norm roundtrip test fails.

  \bigskip

  The g-invariance code indexes $\beta$ by \texttt{k-1} in the loop, so it
  was already correct despite the paper typo.

  \smallskip

  \textbf{Note}: Typo still present in the NeurIPS 2024 camera-ready
  (Eq.~\texttt{eq:new\_coeff\_dihedral} and Algorithm~4, Appendix~E).
\end{frame}

% ============================================================
\begin{frame}{Technical challenges: CG matrices without scipy}
  \textbf{Goal}: find orthogonal $C$ such that
  $(\rho_i \otimes \rho_j)(g) = C \bigl[\bigoplus_\rho \rho(g)\bigr] C^\top$
  for all $g \in D_n$.

  \smallskip

  The g-invariance repo used \texttt{scipy.linalg.schur}. We need pure PyTorch.

  \smallskip

  \textbf{Step 1}: Eigendecompose $Q_a = \operatorname{kron}(\rho_i(a), \rho_j(a))$
  via \texttt{torch.linalg.eig}. This $4\times 4$ orthogonal matrix has
  eigenvalues $e^{\pm i\theta}$ (conjugate pairs $\to$ 2D irreps) or $\pm 1$
  (real $\to$ 1D irreps).

  \smallskip

  \textbf{Step 2 --- 2D blocks}: For each conjugate pair $e^{\pm i\theta}$,
  take $u_1 = \operatorname{Re}(v)$, $u_2 = \operatorname{Im}(v)$.
  In this basis the block is $R(-\theta)$, but our irrep convention is
  $\rho_k(a) = R(+2\pi k/n)$.
  Fix: negate $u_2$ when $\theta > 0$ so the block becomes $R(+|\theta|)$.

  \smallskip

  \textbf{Step 3 --- 1D blocks}: When two 1D irreps share the same eigenvalue
  under $a$ (e.g.\ $\rho_0$ and $\rho_{01}$ both have $+1$), the eigenspace is
  2D and \texttt{eig} won't split them. We restrict
  $Q_x = \operatorname{kron}(\rho_i(x), \rho_j(x))$ to the degenerate
  subspace and diagonalize with \texttt{eigh}; its eigenvalues $\pm 1$ separate
  the two 1D irreps.

  \smallskip

  \textbf{Step 4 --- reflection alignment}: Rotate each 2D subspace by angle
  $\alpha$ so that $C^\top Q_x\, C = \operatorname{diag}(1,-1)$, matching
  the standard irrep form $\rho_k(x) = \bigl[\begin{smallmatrix}1&0\\0&-1\end{smallmatrix}\bigr]$.
\end{frame}

% ============================================================
\begin{frame}{Technical challenges: bispectrum formula}
  Non-commutative bispectrum (Theorem 2.3):
  \[
    \beta_{\rho_1,\rho_2} = \bigl[\mathcal{F}_{\rho_1} \otimes \mathcal{F}_{\rho_2}\bigr]
    \, C \, \bigl[\bigoplus_\rho \mathcal{F}_\rho^\top\bigr] \, C^\top
  \]

  \bigskip

  Getting the transpose/adjoint conventions right for real-valued irreps
  required several iterations. Initial implementation broke $D_n$ invariance
  due to incorrect matrix multiplication order.

  \bigskip

  Verified by testing invariance under all $n$ rotations and reflections
  for $n \in \{3, 4, 5, 7, 8, 16\}$.
\end{frame}

% ============================================================
\begin{frame}{Technical challenges: $O(2)$ indeterminacy in inversion}
  Algorithm 4 recovers $\mathcal{F}(\rho_1)$ from:
  \[
    \frac{\beta_{\rho_0,\rho_1}}{\mathcal{F}(\rho_0)} = \mathcal{F}(\rho_1)^\top \mathcal{F}(\rho_1)
  \]

  This determines $\mathcal{F}(\rho_1)$ only up to $Q \in O(2)$: any $Q \cdot \mathcal{F}(\rho_1)$ also satisfies the equation.

  \bigskip

  \textbf{Consequence}: exact bispectrum roundtrip $\beta \to \text{invert} \to \text{forward} \to \beta' = \beta$ cannot hold for all coefficients. This is a \emph{theoretical} limitation, not a bug.

  \bigskip

  \textbf{What we guarantee}:
  \begin{itemize}
    \item $\beta_{\rho_0,\rho_0}$ and $\beta_{\rho_0,\rho_1}$ roundtrip exactly (they are $O(2)$-invariant)
    \item Frobenius norms $\|\mathcal{F}(\rho_k)\|_F$ match for all $k$
  \end{itemize}
\end{frame}

% ============================================================
\begin{frame}{From group DFT to standard FFT: the derivation}
  The $D_n$ group DFT for 2D irrep $\rho_k$ sums over rotation and reflection elements:
  \[
    \hat{f}(\rho_k) = \sum_{l=0}^{n-1} f_{\text{rot}}(l)\, R\!\left(\tfrac{2\pi kl}{n}\right)^{\!\top}
    + \sum_{l=0}^{n-1} f_{\text{ref}}(l)\, \bigl[R\!\left(\tfrac{2\pi kl}{n}\right) \operatorname{diag}(1,-1)\bigr]^{\!\top}
  \]

  Expand the $2\times 2$ entries, noting $R(\theta)^\top = \bigl[\begin{smallmatrix}\cos\theta & \sin\theta \\ -\sin\theta & \cos\theta\end{smallmatrix}\bigr]$:
  \[
    \hat{f}_{00}^{(k)} = \textstyle\sum_l f_{\text{rot}}(l)\cos\theta_{kl} + \sum_l f_{\text{ref}}(l)\cos\theta_{kl}
  \]

  But $\texttt{fft}(f)[k] = \sum_l f(l)\, e^{-i 2\pi kl/n}$, so
  $\operatorname{Re}(F_{\text{rot}}[k]) = \sum_l f_{\text{rot}}(l) \cos\theta_{kl}$.

  \smallskip

  Therefore all four matrix entries are Re/Im combinations of two standard FFTs:
  \[
    \hat{f}^{(k)} = \begin{bmatrix}
      \operatorname{Re}F_r + \operatorname{Re}F_x & \operatorname{Im}F_r - \operatorname{Im}F_x \\
      -\operatorname{Im}F_r - \operatorname{Im}F_x & \operatorname{Re}F_r - \operatorname{Re}F_x
    \end{bmatrix}
  \]
  where $F_r = \texttt{fft}(f_{\text{rot}})$ and $F_x = \texttt{fft}(f_{\text{ref}})$.
  Two $O(n\log n)$ FFTs replace one $O(n^2)$ matrix--vector product.

  \smallskip

  The inverse is analogous: solve for $F_r[k], F_x[k]$ from the four entries,
  enforce Hermitian symmetry $F[n{-}k] = \overline{F[k]}$, then \texttt{ifft}.
\end{frame}

% ============================================================
\begin{frame}{FFT-based group DFT: 100--275$\times$ speedup}
  Replaced $O(n^2)$ explicit DFT (einsum against precomputed rotation matrices)
  with two calls to \texttt{torch.fft.fft} / \texttt{torch.fft.ifft}.

  \bigskip

  Key insight: each entry of the $2\times 2$ Fourier coefficient $\hat{f}(\rho_k)$
  is a linear combination of $\operatorname{Re}/\operatorname{Im}$ of the standard FFT
  of the rotation and reflection halves of the signal.

  \bigskip

  \begin{table}
    \scriptsize
    \begin{tabular}{lrrrl}
      \toprule
      \textbf{Method} & $n$ & \textbf{Before} & \textbf{After} & \textbf{Speedup} \\
      \midrule
      \texttt{\_group\_dft}   &   64 & 50 ms   & 0.23 ms & $218\times$ \\
      \texttt{\_group\_dft}   &  256 & 76 ms   & 0.28 ms & $275\times$ \\
      \texttt{\_group\_dft}   & 1024 & 89 ms   & 0.44 ms & $205\times$ \\
      \texttt{\_inverse\_dft} &   64 & 2.3 ms  & 0.27 ms & $8.5\times$ \\
      \texttt{\_inverse\_dft} &  256 & 10.5 ms & 0.33 ms & $32\times$ \\
      \texttt{\_inverse\_dft} & 1024 & 71 ms   & 0.56 ms & $127\times$ \\
      \bottomrule
    \end{tabular}
  \end{table}

  Also removes $O(n^2)$ memory from precomputed rotation buffers.
\end{frame}

% ============================================================
% SO2onD2 SECTION
% ============================================================

\begin{frame}{}
  \vfill
  \centering
  {\Large\textbf{New: SO2onD2}}

  \medskip

  {\normalsize Implementing the Selective Disk Bispectrum}

  \smallskip

  {\small Myers \& Miolane, arXiv:2511.19706, 2025}
  \vfill
\end{frame}

% ============================================================
\begin{frame}{SO2onD2: What and why}
  \textbf{Setting}: $f : \mathbb{D} \to \mathbb{R}$ (grayscale image on the unit disk).
  $\mathrm{SO}(2)$ acts by rotation: $(R_\phi f)(r,\theta) = f(r, \theta - \phi)$.

  \medskip

  \textbf{Disk Harmonic Transform} (analog of Fourier/SHT):
  \[
    f(r, \theta) = \textstyle\sum_{n,k} a_{nk} \, \psi_{nk}(r, \theta), \qquad
    \psi_{nk} = c_{nk} \, J_n(\lambda_{nk} r) \, e^{in\theta}
  \]
  \textbf{Equivariance}: $a^{R_\phi f}_{n,k} = e^{in\phi} \cdot a^f_{n,k}$

  \medskip

  \textbf{Selective disk bispectrum} (Def.~4.2) --- $O(m)$ coefficients:
  \[
    b_{0,0,k} = a_{0,1}^2 \cdot a_{0,k}^*, \qquad
    b_{2,n,k} = a_{1,1} \cdot a_{n,1} \cdot a_{n+1,k}^*
  \]
  Phases cancel: $e^{i\phi} \cdot e^{in\phi} \cdot e^{-i(n+1)\phi} = 1$. \textbf{Invertible} (Thm.~4.4).
\end{frame}

% ============================================================
\begin{frame}{SO2onD2: Architecture}
  Follows the same \texttt{nn.Module} contract as \texttt{CnonCn}/\texttt{DnonDn}:
  \begin{itemize}
    \item No trainable parameters (all buffers)
    \item \texttt{forward(f)} takes raw $L \times L$ images, handles DHT internally
    \item \texttt{invert(beta)} recovers image up to rotation
    \item \texttt{output\_size}, \texttt{index\_map} properties
  \end{itemize}

  \bigskip

  \textbf{Key design choices}:
  \begin{enumerate}
    \item No new dependencies --- Bessel functions from scratch in pure torch
    \item No external data files --- all roots \& matrices computed at construction
    \item \texttt{selective=False} raises \texttt{NotImplementedError} (as in \texttt{DnonDn})
    \item Mirrors \texttt{CnonCn} structure: $\mathrm{SO}(2)$ is commutative $\Rightarrow$ scalar triple products, no CG matrices
  \end{enumerate}
\end{frame}

% ============================================================
\begin{frame}{Challenge 1: Bessel functions without scipy}
  Need $J_n(x)$ and its positive zeros $\lambda_{nk}$ where $J_n(\lambda_{nk}) = 0$.

  \bigskip

  \textbf{$J_n(x)$}: Forward recurrence from \texttt{torch.special.bessel\_j0/j1}:
  \[
    J_{k+1}(x) = \frac{2k}{x} J_k(x) - J_{k-1}(x)
  \]
  Stable when $x \geq n$ (our regime: $x = \lambda_{nk} r$, $\lambda_{nk} > n$).

  \bigskip

  \textbf{Root finding}: \textbf{Interlacing property} $\lambda_{n-1,k} < \lambda_{n,k} < \lambda_{n-1,k+1}$ gives guaranteed brackets.

  \begin{enumerate}
    \item $J_0$ roots via McMahon expansion + Newton polish
    \item For $n = 1, 2, \ldots$: bracket each $J_n$ root between consecutive $J_{n-1}$ roots
    \item Newton--bisection hybrid within each bracket (guaranteed convergence)
  \end{enumerate}

  Single-pass algorithm: \texttt{compute\_all\_bessel\_roots(n\_max, k\_max)}.
\end{frame}

% ============================================================
\begin{frame}{Challenge 2: DHT on a discrete pixel grid}
  \textbf{Problem}: The natural complex basis $\{\psi_{nk}, \psi_{-n,k}\}$ is redundant for real signals ($\psi_{-n,k} = \overline{\psi_{n,k}}$). On a discrete grid, the synthesis matrix is \emph{rank-deficient} (condition number $\sim 10^{20}$).

  \bigskip

  \textbf{Solution}: Reformulate using a \textbf{real basis}:
  \begin{align*}
    n = 0: &\quad c_{0k} J_0(\lambda_{0k} r) \hfill \text{(1 column per }(0,k)\text{)} \\
    n > 0: &\quad 2\,c_{nk} J_n(\lambda_{nk} r)\cos(n\theta),\quad -2\,c_{nk} J_n(\lambda_{nk} r)\sin(n\theta) \hfill \text{(2 columns)}
  \end{align*}
  Yields matrix $\Phi \in \mathbb{R}^{p_{\text{disk}} \times d_{\text{real}}}$, well-conditioned ($\kappa \sim 10^3$).

  \bigskip

  \textbf{Analysis}: $\mathbf{x} = \Phi^+ \mathbf{f}$ \quad (pseudoinverse, \texttt{rcond=1e-10})

  \textbf{Synthesis}: $\mathbf{f} = \Phi \, \mathbf{x}$

  \smallskip

  Convert real $\mathbf{x}$ to complex $a_{n,k} = x_{\text{re}} + i\,x_{\text{im}}$ for the bispectrum.
\end{frame}

% ============================================================
\begin{frame}{Challenge 3: Bandlimit convention}
  Paper says $\lambda = \pi L / 2$, but what determines $m$ exactly?

  \bigskip

  Convention from Marshall et al.\ (2023, \texttt{fle-2d}):
  \[
    n_e = \lfloor L^2 \pi / 4 \rfloor \approx \text{number of pixels inside the disk}
  \]
  Sort all $(n, k)$ by ascending $\lambda_{nk}$, keep the first $n_e$.

  \bigskip

  \textbf{Edge case}: if $n_e$ splits a conjugate pair $(n, k)$ and $(-n, k)$, drop the orphan.

  \bigskip

  \begin{table}
    \scriptsize
    \begin{tabular}{lccccc}
      \toprule
      \textbf{Image size} & $L=8$ & $L=16$ & $L=28$ & $L=56$ & $L=112$ \\
      \midrule
      Paper (Table 1) & 27 & 105 & 315 & 1{,}247 & 4{,}957 \\
      Our implementation & 27 & 105 & 315 & --- & --- \\
      \bottomrule
    \end{tabular}
  \end{table}

  \smallskip

  Exact match for all tested sizes. ($L \geq 56$ untested due to construction time.)
\end{frame}

% ============================================================
\begin{frame}{Inversion: bootstrap recovery (Thm.~4.4)}
  Same pattern as \texttt{CnonCn.invert()}, adapted to 2D indexing:

  \medskip

  \begin{enumerate}\setlength{\itemsep}{1pt}
    \item $a_{0,1}$ from $b_{0,0,1} = |a_{0,1}|^3 e^{i\arg(a_{0,1})}$ \quad $\Rightarrow$ \quad cube root
    \item $a_{0,k}$ for $k \geq 2$: \quad $a_{0,k} = \overline{b_{0,0,k} / a_{0,1}^2}$
    \item $|a_{1,1}|$ from $b_{2,0,1} = a_{0,1} |a_{1,1}|^2$ \quad (phase = rotation indeterminacy)
    \item $a_{1,k}$ for $k \geq 2$: \quad $a_{1,k} = \overline{b_{2,0,k} / (a_{1,1} \cdot a_{0,1})}$
    \item Chain for $n \geq 1$: \quad $a_{n+1,k} = \overline{b_{2,n,k} / (a_{1,1} \cdot a_{n,1})}$
  \end{enumerate}

  \medskip

  \textbf{At the coefficient level}: exact (magnitudes match to $10^{-12}$).

  \textbf{Full roundtrip} ($f \to \beta \to \hat{f} \to f'$): limited by DHT discretization;
  ${>}85\%$ of DH coefficients recovered within 20\% relative error.

  \medskip

  \textbf{Assumption}: $a_{n,1} \neq 0$ for all $n$.
  Mild for noisy images; fails for purely radial signals.
\end{frame}

% ============================================================
\begin{frame}[fragile]{SO2onD2: Usage}
\begin{lstlisting}[language=Python]
from bispectrum import SO2onD2

bsp = SO2onD2(L=28)          # 28x28 images
print(bsp.output_size)        # 315

f = torch.randn(4, 28, 28)   # batch of 4 images
beta = bsp(f)                 # (4, 315), complex128
f_rec = bsp.invert(beta)      # (4, 28, 28), float64

# Invariance check
cos_a, sin_a = math.cos(0.7), math.sin(0.7)
theta = torch.tensor([[[cos_a, sin_a, 0],
                        [-sin_a, cos_a, 0]]])
grid = F.affine_grid(theta, [1,1,28,28])
f_rot = F.grid_sample(f[:1].unsqueeze(1), grid)

beta_rot = bsp(f_rot.squeeze(1))
# beta and beta_rot are approximately equal
\end{lstlisting}
\end{frame}

% ============================================================
\begin{frame}{Numerical subtleties encountered}
  \begin{enumerate}
    \item \textbf{Forward recurrence cancellation}: $J_n(x)$ near its own zeros loses $\sim$7 digits for $n \geq 2$.
    Our computed roots are zeros of the \emph{numerically evaluated} $J_n$, not the true $J_n$.
    Everything is self-consistent; displacement $\sim 6 \times 10^{-5}$.

    \smallskip

    \item \textbf{dtype leak}: \texttt{torch.tensor(x)} with Python float infers \texttt{float32}, not \texttt{float64}. Silently halved precision of all root-finding.

    \smallskip

    \item \textbf{Phase precision}: \texttt{torch.exp(1j * n * phi)} creates \texttt{complex64} even when multiplied with \texttt{complex128} tensors. Must explicitly pass \texttt{dtype=complex128}.

    \smallskip

    \item \textbf{Pseudoinverse regularization}: 1--2 basis functions per grid are unresolvable (aliased at the pixel scale). Setting \texttt{rcond=1e-10} in \texttt{pinv} zeros their singular values instead of amplifying noise.
  \end{enumerate}
\end{frame}

% ============================================================
\begin{frame}{Test coverage: 180 tests}
  \begin{columns}[T]
    \begin{column}{0.48\textwidth}
      \textbf{CnonCn} (42 tests)
      \begin{itemize}\scriptsize
        \item Output shape/dtype (6 values of $n$)
        \item Cyclic shift invariance (selective \& full, all shifts)
        \item Formula verification vs manual DFT
        \item Inversion: DFT magnitude roundtrip
        \item Inversion: full bispectrum roundtrip
      \end{itemize}

      \bigskip

      \textbf{DnonDn} (67 tests)
      \begin{itemize}\scriptsize
        \item Rotation invariance (all shifts, 6 values of $n$)
        \item Reflection invariance (6 values of $n$)
        \item Combined rotation + reflection
        \item DFT roundtrip (forward + inverse = identity)
        \item Inversion: $\beta_{\rho_0,*}$ exact roundtrip
        \item Inversion: Fourier Frobenius norm recovery
        \item DFT performance benchmarks ($n$ up to 1024)
      \end{itemize}
    \end{column}
    \begin{column}{0.48\textwidth}
      \textbf{SO2onD2} (45 tests)
      \begin{itemize}\scriptsize
        \item Bessel: $J_n$ vs torch, root accuracy, interlacing
        \item Construction: coeff counts vs paper Table~1
        \item Forward: shape, dtype, analytical \& spatial invariance
        \item DHT roundtrip: bandlimited \& single harmonic
        \item Inversion: coefficient-level, bispectrum roundtrip
      \end{itemize}

      \bigskip

      \textbf{Other} (32 tests)
      \begin{itemize}\scriptsize
        \item SO3onS2: construction, forward, rotation invariance
        \item Rotation utilities, import / public API surface
      \end{itemize}
    \end{column}
  \end{columns}
\end{frame}

% ============================================================
\begin{frame}{Weaknesses \& areas for improvement}
  \begin{enumerate}
    \item \textbf{No GPU testing} --- all tests run on CPU only

    \item \textbf{$D_n$ inversion is partial} --- $O(2)$ indeterminacy is fundamental;
      worth investigating canonical choices of $Q$

    \item \textbf{Full bispectrum not implemented} --- only \texttt{selective=True}
      works for \texttt{DnonDn} and \texttt{SO2onD2}

    \item \textbf{SO2onD2 construction is slow} --- Bessel root finding + pseudoinverse
      in pure Python; $L=28$ takes ${\sim}15$s, $L=64$ impractical

    \item \textbf{Naive DHT} --- $O(L^3)$ pseudoinverse instead of
      $O(L^2 \log L)$ fast DHT from Marshall et al.\ (2023)

    \item \textbf{No G-CNN integration yet} --- standalone module,
      not wired into an equivariant network pipeline
  \end{enumerate}
\end{frame}

% ============================================================
\begin{frame}{Next steps}
  \begin{itemize}
    \item \textbf{SO2onD2: MRA correction} (Proposition 4.6) --- bias term $\delta$ for noisy rotated observations
    \item \textbf{SO2onD2: fast DHT} --- integrate $O(L^2 \log L)$ approximation from Marshall et al.
    \item \textbf{SO2onD2: performance} --- cache Bessel roots, or move to C++/CUDA
    \item Wire into a G-CNN pipeline (replace escnn pooling layer)
    \item Reproduce MNIST/EMNIST/CIFAR-10 experiments from the papers
    \item Explore octahedral / full octahedral groups (Mataigne Appendix F)
    \item Investigate selective bispectrum for $\mathrm{SO}(3)$ on $S^2$ (open problem)
    \item GPU testing and benchmarking
  \end{itemize}
\end{frame}

\end{document}
